%%=============================================================================
%% Voorwoord
%%=============================================================================

\chapter*{\IfLanguageName{dutch}{Woord vooraf}{Preface}}%
\label{ch:voorwoord}

%% TODO:
%% Het voorwoord is het enige deel van de bachelorproef waar je vanuit je
%% eigen standpunt (``ik-vorm'') mag schrijven. Je kan hier bv. motiveren
%% waarom jij het onderwerp wil bespreken.
%% Vergeet ook niet te bedanken wie je geholpen/gesteund/... heeft

Deze bachelorproef is geschreven in het kader van mijn professionele bachelor Toegepaste Informatica aan HOGENT. De weg naar dit onderwerp verliep niet in één keer rechtlijnig. Aanvankelijk werkte ik met een ander onderwerp, maar tijdens een gesprek in het begin van het traject bleek dat dit nog te vaag was afgebakend. In die periode was Mevr. L. Van Steenberghe afwezig en nam lector Lena De Mol tijdelijk over. Dankzij dat gesprek en een korte brainstorm met haar kreeg ik een nieuw, concreter idee dat uiteindelijk leidde tot dit onderzoek.

Die keuze sloot ook aan bij een ervaring uit het vorige semester. Toen werkte ik met mijn team aan een applicatie voor HOGENT. We hebben daar sterke resultaten neergezet, maar bij de uitwerking lag de prioriteit vooral op het opleveren van de applicatie zelf en minder op toegankelijkheid. Dat bleef bij mij hangen. Sindsdien stelde ik mij steeds vaker de vraag waarom digitale toegankelijkheid zo weinig prioriteit krijgt, terwijl net dat aspect bepaalt of iedereen volwaardig kan deelnemen. Vanuit die overtuiging groeide mijn motivatie om digitale toegankelijkheid aan HOGENT te onderzoeken vanuit het perspectief van studenten met een beperking.

Het onderwerp bleek maatschappelijk relevant en tegelijk technisch uitdagend. Toegankelijkheid raakt aan wetgeving, ontwerpkeuzes, gebruiksvriendelijkheid en inclusie. Dat maakte het afbakenen van de scope niet altijd eenvoudig. Tijdens dit traject heb ik geleerd om kritischer naar digitale omgevingen te kijken, onderzoeksmatig te werken en theorie te vertalen naar concrete aanbevelingen voor de praktijk.

Graag bedank ik in de eerste plaats mijn promotor mevrouw L. Van Steenberghe, voor de waardevolle feedback, de gerichte begeleiding en het vertrouwen tijdens dit traject. Daarnaast wil ik mevrouw Lena De Mol expliciet bedanken voor haar begeleiding bij de opstart en voor de richtinggevende input die mee aan de basis lag van dit onderwerp. Ook bedank ik mijn copromotor Dhr. A. De Waele, voor de inhoudelijke ondersteuning en de praktische inzichten vanuit het werkveld.

